\section{Related Work}
\label{sec:related}

\emph{Speculative decoding.} Leviathan et al.~\cite{leviathan2023} and Chen et al.~\cite{chen2023} introduced draft-then-verify decoding with verification rules that preserve the target model's output distribution; Xia et al.~\cite{xia2024survey} survey the many variants that followed. One line of work raises the number of tokens accepted per target pass. SpecInfer organizes the candidates of small speculative models into a token tree that the target verifies in parallel~\cite{specinfer2024}. Sequoia constructs token trees by dynamic programming and adapts their size and depth to the hardware~\cite{sequoia2024}. BiLD lets a small model generate most tokens and invokes the large model only when needed~\cite{bild2023}. Medusa~\cite{medusa2024} and EAGLE~\cite{eagle2024} remove the separate draft model by adding decoding heads to the target or drafting at the feature level. Lookahead decoding~\cite{lookahead2024} parallelizes decoding without any draft model; despite its name, it is unrelated to the draft-ahead lookahead studied here. These techniques assume that both models share one machine. We use strict greedy verification with an independent draft model on the device. \cape{} changes only the drafting schedule around the network round trip, which is complementary to techniques that raise acceptance.

\emph{Adaptive draft length.} AdaEDL~\cite{adaedl2024} stops drafting when an entropy-based lower bound on the acceptance probability becomes small, and SpecDec++~\cite{specdecpp2025} trains an acceptance-prediction head to decide when to stop. These methods decide how many tokens to propose in a co-located setting where drafting and verification alternate. \cape{} instead decides how many tokens to draft on and off the critical path of an edge--cloud pipeline. It needs neither draft probabilities nor training, and it prices the device energy of speculative work explicitly. Our comparison with AdaEDL (Section~\ref{sec:main}) shows why the distinction matters: when verification cost is nearly independent of the draft length, stopping early saves little.

\emph{Parallel drafting and verification.} PEARL~\cite{pearl2025} combines pre-verification and post-verification so that drafting and verification wait less for each other, and it adapts the draft length. Overlapping the two models is therefore not new. \cape{} applies the overlap across a network and contributes a model that prices both the critical path and the energy of speculative work, for a device that pays for every discarded token.

\emph{On-device and edge--cloud speculative decoding.} EdgeLLM accelerates LLM inference on mobile devices with speculative decoding~\cite{edgellm2025}, keeping both models on the device. Several systems instead split the models across machines. In SLED, lightweight edge devices draft with their own draft models and a shared edge server verifies their requests in batches~\cite{sled2025}. SpecEdge drafts on consumer-grade edge GPUs, overlaps edge drafting with server verification, and interleaves the requests of many users to raise server throughput~\cite{specedge2025}. DSD studies distributed speculative decoding with a discrete-event simulator and an adaptive window policy~\cite{dsd2025}. PicoSpec runs edge and cloud models in an asynchronous pipeline and compresses the transmitted vocabulary distributions~\cite{picospec2026}. WISP suppresses wasted drafting and verification interference with dynamic drafting and SLO-aware batching~\cite{wisp2026}. Lyu et al.~\cite{distance2026} show analytically that pipelining makes distributed speculative decoding competitive mainly when the round trip is short relative to edge drafting, and that its main advantage is server capacity rather than single-request latency. These systems mostly target server-side throughput, multi-user scheduling or communication volume. \cape{} targets the critical path and the energy of a single device: it runs entirely on the device, treats the verifier as a black box, and decides per round how much drafting to expose and how much to hide, which is exactly the trade-off that determines whether pipelining pays off.

\emph{Energy.} Samsi et al.~\cite{samsi2023} benchmark the energy cost of LLM inference on data-center GPUs. Dutta et al.~\cite{energy2026} show that the speedup of speculative decoding does not translate proportionally into energy savings, which motivates measuring energy directly rather than treating speed as its proxy. We measure the board energy of the edge device and use a fitted energy model inside the controller.
